%% kapa: exam=1 round=28 year=2017 period=2 form=A subject=appraisal_law subject_ko=감정평가관계법규 seq=1 number=1
%% source: Q-Net 감정평가사 기출문제 (cst00309) artlSeq=5199603 file=extracted/1008/5199603/2017년 제28회 감정평가사 1차시험 문제지/2교시 A형.pdf
%% conversion: kordoc 4.24.51 -> tools/gen1.py (auto); review=proofread-sonnet; flags=none
\begin{question}{1}
국토의 계획 및 이용에 관한 법령상 도시ㆍ군관리계획의 입안 등에 관한 설명으로 옳지 않은 것은?
\begin{choices}
  \item 주민은 기반시설의 개량에 관한 사항에 대하여 도시ㆍ군관리계획의 입안을 제안할 수 있다.
  \item 도시ㆍ군관리계획의 입안을 제안받은 자는 제안자와 협의하여 제안된 도시ㆍ군관리계획의 입안 및 결정에 필요한 비용의 전부 또는 일부를 제안자에게 부담시킬 수 있다.
  \item 지구단위계획구역의 지정에 관한 사항에 대하여 도시ㆍ군관리계획의 입안을 제안하려는 자는 국ㆍ공유지를 제외한 대상토지면적의 3분의 2 이상의 토지소유자의 동의를 받아야 한다.
  \item 도시ㆍ군관리계획으로 입안하려는 지구단위계획구역이 상업지역에 위치하는 경우에는 재해취약성분석을 실시하여야 한다.
  \item 도시지역의 축소에 따른 지구단위계획구역의 변경에 대한 도시ㆍ군관리계획을 입안할 때에는 주민의 의견청취가 요구되지 아니한다.
\end{choices}
\end{question}
\answer{4}
